\chapter{Concluzii}

În cadrul acestei lucrări, obiectivele propuse au fost îndeplinite cu succes prin dezvoltarea unei platforme web complexe de programare competitivă. Proiectul a demonstrat viabilitatea integrării unei arhitecturi client-server utilizând tehnologii moderne: React și TypeScript pentru interfață, respectiv Spring Boot și PostgreSQL pentru logica de backend și persistența datelor. Integrarea cu succes a motorului de execuție Judge0 și a modelelor generative de limbaj (Qwen2.5-coder) a permis automatizarea fluxurilor de evaluarea a codului și de generare a conținutului educațional. 

Consider că rezultatele obținute sunt relevante în contextul educațional actual din domeniul IT. Învățarea algoritmilor este adesea percepută ca fiind rigidă și monotonă, iar prin adăugarea unei componente solide de gamificare procesul este transformat într-o experiență interactivă și mai atractivă. 

Deși aplicația este complet funcțională în stadiul curent, flexibilitatea arhitecturii permite extinderi viitoare de interes academic și tehnic:

\begin{itemize}
    \item Extinderea suportului pentru LLM: migrarea către modele generative mai mari găzduite în cloud sau fine-tuning-ul unui model local dedicat exclusiv generării de probleme algoritmice complexe.
    \item Module de competiție în timp real: implementarea unui sistem de meciuri directe între utilizatori.
    \item Dezvoltarea unei aplicații mobile: transpunerea componentelor de gamificare (în special mecanismul de chestionarul zilei) într-o aplicație mobilă pentru a facilita accesul rapid al utilizatorilor și a crește rata de retenție.
\end{itemize}